\documentclass[10pt,conference,compsocconf]{IEEEtran}

\usepackage{hyperref}
\usepackage{graphicx}	% For figure environment


\begin{document}
\title{Predicting Coronary Heart Disease from Health Survey Data
}

\author{
  Aurore Zhang, Jiayi Fang and Vincent Coulombeau \\
  \textit{Computer Science students, CS-433, EPFL, Switzerland}
}

\maketitle

\begin{abstract}
  Cardiovascular diseases are among the leading causes of death worldwide,
  and early identification of people at risk can support prevention. In this
  project, we use the 2015 BRFSS health survey to predict whether a respondent
  has coronary heart disease or has had a heart attack.
\end{abstract}

\section{Introduction}

According to the World Health Organization, Cardiovascular Diseases (CVD),
such as heart attacks, are becoming one of the leading causes of death
globally. Adults are living longer, and diseases of the heart and
circulatory vessels are prevalent in this growing population of older
adults. The rise of new technologies such as Machine Learning algorithms
can help with the early detection and prevention of developing CVDs.
We shall study in this report a robust predictor of whether a person has
coronary heart disease or has suffered a heart attack using the BRFSS
2015 health survey. It is a binary classification problem where:
\begin{itemize}
  \item Negative class: -1
  \item Positive class: +1
  \item Input: 321 survey features
  \item Training examples: approximately 328,135
  \item Positive class: approximately 8.83\%
\end{itemize}
The complete implementation is launched by run.py.

\section{Data Preprocessing}
\label{sec:structure-paper}

Before training, the raw BRFSS CSVs are loaded and cached, then
split into train/development/holdout while preserving class balance.
The pipeline first removes survey-administration columns and decodes
BRFSS sentinel codes such as “don’t know/refused/none” into missing
values or zero, because these codes are not real medical measurements
and would otherwise distort the model. Missing values are then
imputed using the training-set median, and the data are standardized
so features on different scales contribute comparably. Low-cardinality
categorical variables are one-hot encoded, composite banded codes are
split into count + band features, and missingness indicators and
selected interaction terms are added to capture informative response
patterns and simple nonlinear effects.
Finally, all these transformations are fitted only on the training
data and then applied unchanged to validation/test data, preventing
data leakage and ensuring the model sees a clean, consistent
representation of the raw survey responses.
These transformations are applied optionally depending on the
arguments passed when running the script, this will allow us to
evaluate the relevance and the efficiency of these techniques when
evaluating the predictions of the test set.

\section{Models and Experiments}
\label{sec:models}

We use ridge regression from \texttt{implementations.py} with the input features prepared using the preprocessing steps described in Section~\ref{sec:structure-paper}. Since only $8.83\%$ of
rows are positive, an unweighted model tends to predict the negative class.
We give positive examples a weight $w_{+}$ and minimise
\begin{equation}
  \label{eq:weighted-ridge}
  \frac{1}{2N}\sum_{i=1}^{N}s_i(y_i-x_i^\top w)^2
  +\lambda\lVert w\rVert^2,
\end{equation}
where $s_i=w_{+}$ for positive examples and $s_i=1$ otherwise. We implement
this by multiplying each row and target by $\sqrt{s_i}$ before calling
\texttt{ridge\_regression}. We tune the decision threshold using development
labels over $100$ score quantiles. For refitted models, we transfer the
fraction of rows predicted positive, rather than the score threshold, to
account for changes in score scale. Class weighting changes the fitted
model, while threshold tuning changes which scores count as positive.

We split raw rows by class into $60/20/20$ training, validation and test
sets before preprocessing. All preprocessing steps are fitted only on the
training data. We select model settings using either the validation set or stratified 5-fold cross-validation on the combined training and validation sets (80\% of the data). During cross-validation, each fold’s decision threshold is tuned on an inner validation split of that fold’s training data. We evaluate the final model once on the held-out test set, then refit on all 328,135 labelled rows to generate submission predictions.

 
We run four experiments:
(i) compare \texttt{standardize} ($322$ columns), \texttt{all} ($217$) and
\texttt{clean} ($789$) on the same splits;
(ii) test $\lambda\in\{10^{-5},\ldots,1\}$ and $w_{+}\in\{1,3,5,7\}$ in
$72$ fits, including nearly unregularised and unweighted models, with a
majority-class predictor as the baseline;
(iii) add preprocessing stages one at a time: standardisation, removal of
administrative columns, sentinel decoding, missingness indicators and the
remaining feature construction;
(iv) use paired $5$-fold F1 differences to compare clipping, one-hot width,
interaction count, excluding binary columns, curated nominal encoding and
bagging. We also test regularised logistic regression on the same matrix,
with and without class weighting, using
$\gamma\in\{0.05,0.2,1,4,16\}$ and up to $4{,}000$ iterations.

\section{Results}
\label{sec:results}

\textit{Holdout performance.}
For the single-split selection (\texttt{clean}, $\lambda=10^{-3}$,
$w_{+}=7$), ridge achieved F1 $0.4311$ on the sealed holdout
(Table~\ref{tab:holdout}). It recovered $53.93\%$ of positive cases,
whereas the majority baseline found none despite its higher accuracy.

\begin{table}[tbp]
  \centering
  \caption{Positive-class performance on the sealed holdout.}
  \label{tab:holdout}
  \small
  \setlength{\tabcolsep}{4pt}
  \begin{tabular}{lrrrr}
    \hline
    Model & Accuracy & Precision & Recall & F1 \\
    \hline
    Always $-1$ & 0.9117 & 0.0000 & 0.0000 & 0.0000 \\
    Selected ridge & 0.8743 & 0.3591 & 0.5393 & 0.4311 \\
    \hline
  \end{tabular}
\end{table}

\textit{Preprocessing ablation.}
Development F1 rose from $0.4089$ to $0.4318$ across the cumulative
stages in Table~\ref{tab:ablation}. Sentinel decoding produced the
largest individual increase ($+0.0109$); removing administrative
columns alone slightly reduced F1.

\begin{table}[tbp]
  \centering
  \caption{Cumulative preprocessing ablation on development data.
  Deltas are relative to the preceding row.}
  \label{tab:ablation}
  \small
  \setlength{\tabcolsep}{3pt}
  \begin{tabular}{lrr}
    \hline
    Stage & F1 & $\Delta$F1 \\
    \hline
    E0: median imputation + scaling & 0.4089 & --- \\
    E1: + remove administrative columns & 0.4079 & $-0.0010$ \\
    E2: + sentinel decoding & 0.4188 & $+0.0109$ \\
    E3: + missingness indicators & 0.4225 & $+0.0037$ \\
    E4: + remaining feature construction & 0.4318 & $+0.0093$ \\
    \hline
  \end{tabular}
\end{table}

\textit{Parameter and model comparisons.}
Joint cross-validated tuning gave a best mean F1 of $0.4384$ at
$\lambda=10^{-4}$ and $w_{+}=3$. Nearby settings formed a broad plateau:
the top five combinations differed by at most $0.0004$, below the
fold standard deviation of about $0.0034$.
Paired five-fold comparisons favoured clipping ($t=2.79$), but not
increasing interactions from $20$ to $30$ ($\Delta\mathrm{F1}=-0.0006$,
$t=-0.93$, one fold won) or one-hot width from $15$ to $40$
($-0.0001$, $t=-0.39$, two folds won). Excluding binary columns,
curated nominal encoding and bagging also gave no consistent gain.
On the same development split and feature matrix, ridge reached
$0.4389$, versus $0.4317$ for the best measured logistic run
($\gamma=0.2$, $4{,}000$ iterations), a gap of $0.0072$.

\textit{Competition results.}
The recorded competition F1 increased from $0.419$ for the first
submission to $0.448$ for the improved full-data refit. This external
test score and the local holdout score measure different fitted models
on different rows and should not be directly equated.


\section{Discussion and Summary}
\label{sec:discussion}

The results suggest that most improvements came from correcting the
values within columns rather than adding or removing features. With about
$650$ training rows per feature, adding more features did not seem to help.
Removing administrative columns and adding more one-hot features or
interactions made little difference. Sentinel decoding was more useful
because survey response codes should not be treated as numerical values.
For example, a refusal to answer the general-health question is coded as
$9$. Without decoding, the model treats this as worse than ``poor'' health.
These codes also affect the mean and standard deviation used for scaling.

Choosing model settings was less straightforward. Many differences between
settings were smaller than the variation across cross-validation folds.
A setting that looked better on one validation split could perform worse
on every fold in a paired comparison. We therefore used cross-validation
to check these choices and rejected several changes that did not improve F1 consistently. Logistic regression also did not
produce a better ranking, despite avoiding a drawback of squared error:
it can penalise predictions that have the correct sign but go beyond the
$\pm 1$ targets. However, this comparison has a limitation: ridge was solved
directly, while logistic regression used gradient descent and may not have
fully converged.

Our approach has three main limitations. First, the ridge regression
also penalises the constant term, which may affect predictions when the classes
are highly imbalanced. Second, we only split composite survey codes when
all decoded values fall between $101$ and $399$. Questions with extra codes,
such as ``none'' or a fourth band, are therefore not fully decoded.
Third, the test score measures the model before the final refit. The
submitted model uses about $67\%$ more training rows, so its performance
may differ. Overall, giving positive cases more weight and using cleaned survey values, missing-value indicators and one-hot encoded categories helped the ridge model identify more positive cases, but reduced its overall accuracy. This trade-off fits our aim of identifying positive cases,
since the majority-class baseline achieves high accuracy while finding
none.


\end{document}